% ===== PRÉAMBULE CANONIQUE — à coller VERBATIM en tête de CHAQUE .tex =====
\documentclass[11pt,a4paper]{article}
\usepackage[utf8]{inputenc}\usepackage[T1]{fontenc}\usepackage{lmodern}
\usepackage[french]{babel}
\usepackage{amsmath,amssymb,mathtools}\usepackage{icomma}
\usepackage{siunitx}\sisetup{locale=FR}  % \qty{}{}, \si{}, \ang{}
\usepackage{xcolor}\usepackage{colortbl}
\usepackage{tikz}\usepackage{pgfplots}\pgfplotsset{compat=1.18}
\usepackage[siunitx,europeanresistors]{circuitikz} % circuits (élec.)
\usepackage[most]{tcolorbox}\usepackage{enumitem}\usepackage{array,booktabs}
\usepackage[a4paper,margin=2.2cm]{geometry}
\usetikzlibrary{arrows.meta,calc,angles,quotes,positioning,patterns,%
  backgrounds,shapes.geometric,decorations.pathmorphing,decorations.markings,intersections}
\definecolor{primary}{RGB}{222,49,99}\definecolor{primarydark}{RGB}{150,22,55}
\definecolor{defcol}{RGB}{0,114,178}\definecolor{methcol}{RGB}{52,92,148}
\definecolor{excol}{RGB}{0,135,90}\definecolor{attncol}{RGB}{200,40,55}
\tcbset{boxbase/.style={enhanced,breakable,boxrule=0.9pt,arc=2.5pt,
  left=3mm,right=3mm,top=2.3mm,bottom=2.3mm,fonttitle=\bfseries,coltitle=white}}
% --- boîtes TUTORIEL (noms EXACTS, ne pas renommer) ---
\newtcolorbox{cle}[1][]{boxbase,colback=defcol!6,colframe=defcol,
  title={\textsc{À retenir}\ifx\relax#1\relax\relax\else\textmd{ : \itshape #1}\fi}}
\newtcolorbox{methode}[1][]{boxbase,colback=methcol!7,colframe=methcol,sharp corners,
  title={\textsc{Méthode}\ifx\relax#1\relax\relax\else\textmd{ --- \itshape #1}\fi}}
\newtcolorbox{exemplebox}[1][]{boxbase,colback=excol!5,colframe=excol,
  title={\textsc{Exemple}\ifx\relax#1\relax\relax\else\textmd{ : \itshape #1}\fi}}
\newtcolorbox{piege}[1][]{boxbase,colback=attncol!6,colframe=attncol,
  title={\textsc{Piège}\ifx\relax#1\relax\relax\else\textmd{ : \itshape #1}\fi}}
\newtcolorbox{remarque}[1][]{boxbase,colback=black!4,colframe=black!45,
  title={\textsc{Remarque}\ifx\relax#1\relax\relax\else\textmd{ : \itshape #1}\fi}}
% --- boîtes EXERCICE / CORRIGÉ (noms EXACTS, auto-comptées) ---
\newtcolorbox[auto counter]{exo}[1][]{boxbase,colback=primary!4,colframe=primary,
  title={\textsc{Exercice}~\thetcbcounter\ifx\relax#1\relax\relax\else\textmd{ --- \itshape #1}\fi}}
\newtcolorbox[auto counter]{corrige}[1][]{boxbase,colback=excol!4,colframe=excol,
  title={\textsc{Corrigé}~\thetcbcounter\ifx\relax#1\relax\relax\else\textmd{ --- \itshape #1}\fi}}
\newcommand{\vv}[1]{\vec{#1}}\newcommand{\ds}{\displaystyle}
\newcommand{\R}{\mathbb{R}}\newcommand{\N}{\mathbb{N}}\newcommand{\Z}{\mathbb{Z}}
% =========================================================================
\begin{document}
% THEME_TITLE: Le mouvement circulaire uniforme (MCU)
\sisetup{per-mode=symbol}

\noindent Un objet décrit un \textbf{MCU} si sa trajectoire est un cercle (rayon $R$) parcouru à
\emph{vitesse constante en norme}. Bien que la vitesse (norme) soit constante, le mouvement est
\textbf{accéléré} : la \emph{direction} du vecteur vitesse change sans cesse.

\begin{cle}[les grandeurs du MCU]
\begin{itemize}[leftmargin=1.5em,itemsep=2pt]
  \item \textbf{Période} $T$ (durée d'un tour, en \si{\second}) et \textbf{fréquence} $f=\dfrac{1}{T}$ (en \si{\hertz}).
  \item \textbf{Vitesse angulaire} : $\ \omega=\dfrac{2\pi}{T}$ (en \si{\radian\per\second}).
  \item \textbf{Vitesse linéaire} (tangentielle) : $\ \boxed{v=\dfrac{2\pi R}{T}=\omega R}$ (en \si{\metre\per\second}).
  \item \textbf{Accélération et force centripètes} (vers le centre) :
        $\ \boxed{a_c=\dfrac{v^2}{R}}\qquad \boxed{F_c=m\,\dfrac{v^2}{R}=m\,a_c}$.
\end{itemize}
\end{cle}

\begin{center}
\begin{tikzpicture}[>=Stealth,scale=1]
  \def\R{2.2}
  \draw[primary,line width=1pt] (0,0) circle (\R);
  \fill (0,0) circle (1.3pt) node[below left]{$O$};
  \foreach \a in {0,45,...,315}{ \draw[->,defcol,line width=0.8pt] (\a:\R) -- ++(\a+90:0.8); }
  \coordinate (A) at (35:\R);
  \draw[black!55] (0,0) -- (A); \node[black!70,font=\small] at (0.72,0.7){$R$};
  \fill[primarydark] (A) circle (2.3pt) node[above right]{$A$};
  \draw[->,line width=1.1pt,attncol] (A) -- ($(A)!0.42!(0,0)$) node[midway,below right,font=\small]{$\vv F_c$};
  \node[defcol,font=\small,above left] at (125:\R){$\vv v$};
  \node[font=\small,attncol] at (0,-0.55){$a_t=0$};
\end{tikzpicture}
\end{center}

\begin{methode}[résoudre un MCU]
\begin{enumerate}[label=\arabic*),leftmargin=1.7em,topsep=2pt,itemsep=1.5pt]
  \item De $T$ (ou $f$) : $\omega=\tfrac{2\pi}{T}$ et $v=\tfrac{2\pi R}{T}=\omega R$.
  \item Accélération \emph{tangentielle} nulle ($a_t=0$) ; accélération \emph{centripète} $a_c=\tfrac{v^2}{R}$ (vers $O$).
  \item Force centripète $F_c=m\,a_c=m\tfrac{v^2}{R}$ : c'est elle qui « courbe » la trajectoire (tension du fil, frottement, gravité\dots).
\end{enumerate}
\end{methode}

\begin{exemplebox}[une pierre au bout d'un fil]
$m=\qty{0,2}{\kilogram}$, $R=\qty{0,5}{\metre}$, un tour en $T=\qty{1}{\second}$ :
$v=\tfrac{2\pi\times\num{0,5}}{1}\approx\qty{3,14}{\metre\per\second}$,
$\omega=\tfrac{2\pi}{1}\approx\qty{6,28}{\radian\per\second}$,
$a_c=\tfrac{v^2}{R}\approx\qty{19,7}{\metre\per\second\squared}$,
$F_c=m\,a_c\approx\qty{3,9}{\newton}$ (fournie par la tension du fil).
\end{exemplebox}

\begin{piege}[les pièges du MCU]
\begin{itemize}[leftmargin=1.4em,topsep=2pt,itemsep=1.5pt]
  \item Le MCU est \textbf{accéléré} même si $v$ (norme) est constante : la \emph{direction} de $\vv v$ change ($a_t=0$ mais $a_c\ne 0$).
  \item $F_c\propto v^2$ : pour \textbf{doubler} $v$, il faut \textbf{quadrupler} $F_c$ (et non le doubler).
  \item Si le fil casse, l'objet part en \textbf{ligne droite tangente} (inertie), pas « vers l'extérieur ». La force centrifuge est \emph{fictive}.
  \item $v=\omega R$ : à même $\omega$, un point du \emph{bord} va plus vite qu'un point proche du centre.
\end{itemize}
\end{piege}
\end{document}
